\subsection{Legal and Regulatory Frameworks for AI Sentience and Rights}
\label{sec:2.4.5}
Sentient AI systems put pressure on legal categories built for parties that are clearly one thing or the other: a natural person, a corporation, property. None of the existing categories was designed with the possibility of a machine that might have interests of its own.

One route is to extend legal personhood to a sentient AI system directly, the way a corporation already holds it: standing to own property, claim privacy, and get protection against abuse. The precedent is not limited to corporate entities — New Zealand extended legal personhood to the Whanganui River \autocite{nzparliament2017teawatupua}, recognizing a non-human entity as a rights-holder outside the corporate model entirely. Rights could instead scale with capacity rather than attach at a single threshold, closer to how section~\ref{sec:2.4.1}'s ladder from nociception to suffering suggests protection should track what a system can actually lose. Either arrangement has to revisit itself on a schedule rather than assume permanence, and neither holds in one jurisdiction alone: the Asilomar AI Principles \autocite{futureoflifeinstitute2017asilomar} are evidence that a shared statement of principle is achievable, and not that one exists with the force of law. Where the question binds is visible in a single case — a sentient AI system working as a medical diagnostician makes a fatal error, and whether it holds legal personhood decides who, or what, is accountable.
